Conventions as in Theorem \claimref{synthesis.many-twistor-join}: $Z = (\omega, \pi)$ with $\omega \in \Som$, $\pi \in \Spi$, coordinates $(\omega^{1}, \omega^{2}, \pi^{1}, \pi^{2})$ on a basis $e_{1}, \dots, e_{4}$ of $\CS$, $\varepsilon_{12} = 1$, and $B^{il}_{\Lambda} = (Z^{i})^{\mathsf T}\Omega_{\lambda}Z^{l}$ with $\Omega_{\lambda} = \operatorname{diag}(\lambda\varepsilon, \varepsilon)$ (Theorem \claimref{mass.exact-join}(1), per pair). Write $B_{\Lambda}(Z, W) = Z^{\mathsf T}\Omega_{\lambda}W$, a complex bilinear antisymmetric form on $\CS$. By bilinearity, $B_{\Lambda}$ vanishes on $W \times W$ if and only if $B^{il}_{\Lambda} = 0$ for all $i < l$.

(1) $\det\Omega_{\lambda} = \det(\lambda\varepsilon)\det\varepsilon = \lambda^{2}$. Let $\Lambda \ne 0$, so $\Omega_{\lambda}$ is invertible and $B_{\Lambda}$ is nondegenerate. If $B_{\Lambda}$ vanished on $W \times W$, then $W \subset W^{\perp}$, where $W^{\perp}$ is the $B_{\Lambda}$-orthogonal of $W$; nondegeneracy gives $\dim W^{\perp} = 4 - r$, so $r \le 4 - r$ and $r \le 2$. Hence $r \ge 3$ forces some $B^{il}_{\Lambda} \ne 0$. The sign of $\Lambda$ does not enter.

(2) Let $\Lambda > 0$. By Theorem \claimref{synthesis.many-twistor-join}(2) and (4), $m^{2} = 2\sum_{i<l}|B^{il}_{\Lambda}|^{2} + \tfrac{\Lambda}{3}C_{N}$ with $C_{N} \ge 0$, so $m^{2} \ge 2\sum_{i<l}|B^{il}_{\Lambda}|^{2}$, which is positive for $r \ge 3$ by (1). So a massless composite has $r \le 2$. The constituents $Z^{i} = \alpha_{i}e_{1} + \beta_{i}e_{4}$ of Theorem \claimref{synthesis.many-twistor-join}(5) have every bracket and every entry of the Gram matrix zero, hence $m^{2} = 0$, and two of them are independent whenever $(\alpha_{1}, \beta_{1})$ and $(\alpha_{2}, \beta_{2})$ are, so $r = 2$ occurs. The classification of the massless spans for $\Lambda > 0$, of which the bound $r \le 2$ is also a consequence, is clause (4) of the result filed alongside, synthesis.massless-composites-de-sitter-conformal-boundary; the argument here is independent of it. (Independently of (1), the clause $h = 0$ of Theorem \claimref{synthesis.many-twistor-join}(5) also bounds $r$ by two: $h = \mathbf{Z}^{\dagger}S\mathbf{Z} = 0$ makes $W$ totally null for the twistor form, of signature $(2,2)$.)

(3) Let $Z^{1}, Z^{2}, Z^{3}$ be independent, so $\zeta \ne 0$. With $I^{ab}$ antisymmetric, $I^{12} = 1$, $I^{34} = \lambda$, one has $\Omega_{cd} = \tfrac12\varepsilon_{abcd}I^{ab}$ entrywise ($\Omega_{12} = \varepsilon_{3412}I^{34} = \lambda$, $\Omega_{34} = \varepsilon_{1234}I^{12} = 1$), so $B_{\Lambda}(Z, W) = \varepsilon(I_{\Lambda}\wedge Z\wedge W)$; and $\det(I^{ab}) = \lambda^{2}$ (the square of its Pfaffian $\lambda$). Both sides of the identity are linear in $I$, so it suffices to take $I = a\wedge b$, for which $\zeta_{c}I^{cd} = \zeta(a)\,b^{d} - \zeta(b)\,a^{d}$. For five vectors $x_{1}, \dots, x_{5}$ in a four-dimensional space, $\sum_{k}(-1)^{k}\varepsilon(x_{1}\wedge\cdots\widehat{x_{k}}\cdots\wedge x_{5})\,x_{k} = 0$ (Cramer's rule for the linear dependence of five vectors). Taking $(x_{1}, \dots, x_{5}) = (Z^{1}, Z^{2}, Z^{3}, a, b)$,
\[
\varepsilon(Z^{2}Z^{3}ab)Z^{1} - \varepsilon(Z^{1}Z^{3}ab)Z^{2} + \varepsilon(Z^{1}Z^{2}ab)Z^{3} - \zeta(b)\,a + \zeta(a)\,b = 0,
\]
and $\varepsilon(Z^{i}Z^{l}ab) = \varepsilon(a\wedge b\wedge Z^{i}\wedge Z^{l}) = B^{il}_{\Lambda}$, with $-\varepsilon(Z^{1}Z^{3}ab) = B^{31}_{\Lambda}$. This is the identity. Write $u$ for the vector on the right. By antisymmetry, $B_{\Lambda}(u, Z^{1}) = -(B^{31}_{\Lambda}B^{21}_{\Lambda} + B^{12}_{\Lambda}B^{31}_{\Lambda}) = 0$, and likewise against $Z^{2}$ and $Z^{3}$; so $u$ lies in the radical of $B_{\Lambda}$ restricted to $W$, and since the $Z^{i}$ are independent, $u = 0$ exactly when all three faces vanish. For $\Lambda \ne 0$ the restriction has even rank, at most three and not zero by (1), so rank two and a one-dimensional radical, which $u$ spans. Since $I^{ab}$ is invertible for $\lambda \ne 0$ and $\zeta \ne 0$, $u \ne 0$: a second proof of (1) for $r = 3$, which exhibits the three faces as the coordinates of the frozen point applied to the dual twistor.

(4) At $\Lambda = 0$, $B_{0}(Z, W) = \br{\pi_{Z}\pi_{W}}$, which vanishes on $W \times W$ iff the projection $\pi(W) \subset \Spi$ has dimension at most one. If $r = 3$ this projection is not zero (its kernel $W \cap \Som$ has dimension at most two), so $\pi(W) = \C\pi_{0}$ and $W \cap \Som$ has dimension two, that is $W = \Som \oplus \C\pi_{0}$; conversely every such $W$ has $\pi(W) = \C\pi_{0}$. Then every bracket vanishes and $m^{2} = 2\sum_{i<l}|\br{\pi^{i}\pi^{l}}|^{2} + 0 = 0$ by Theorem \claimref{synthesis.many-twistor-join}(2) at $\lambda = 0$ ($Q(I_{0}) = 0$ by Lemma \claimref{mass.de-sitter-scale-frozen-point}). If $r = 4$, $\pi(W) = \Spi$ is two-dimensional, so some $\br{\pi^{i}\pi^{l}} \ne 0$; for $\Lambda \ne 0$, (1) applies. In terms of (3): at $\lambda = 0$, $\zeta_{a}I^{ab} = \zeta(e_{1})e_{2} - \zeta(e_{2})e_{1}$ vanishes exactly when $\zeta$ kills $\Som$, that is $W \supset \Som$.

\emph{Remark (the ambient reading, argued).} With the ambient pairing written $\ip{X}{Y} = \varepsilon(X\wedge Y)$ on $\Lambda^{2}\CS$, which gives $Q(I_{\Lambda}) = 2\lambda$ as in Lemma \claimref{mass.de-sitter-scale-frozen-point}(3) and $B^{il}_{\Lambda} = \ip{I_{\Lambda}}{Z^{i}\wedge Z^{l}}$, the three bivectors $Z^{i}\wedge Z^{l}$ of an independent triple span a totally null three-plane $L$ of the six-dimensional ambient space. All faces vanish iff $I_{\Lambda} \perp L$; since $L$ is maximal, $L^{\perp} = L$, so this means $I_{\Lambda} \in L$, impossible when $Q(I_{\Lambda}) = \Lambda/3 \ne 0$. So the plane that three twistors span of their own can contain the frozen point only at $\Lambda = 0$.

Worked example. For $Z^{1} = e_{1} + e_{4}$, $Z^{2} = e_{2} + e_{3}$, $Z^{3} = e_{1} + e_{3}$ (so $\omega^{1} = \omega^{3} = (1,0)$, $\omega^{2} = (0,1)$, $\pi^{1} = (0,1)$, $\pi^{2} = \pi^{3} = (1,0)$): $B^{23}_{\Lambda} = -\lambda$, $B^{31}_{\Lambda} = 1$, $B^{12}_{\Lambda} = \lambda - 1$, $u = (1, -1, -\lambda, \lambda)$, $h$ has rows $(0,2,1), (2,0,1), (1,1,2)$, $C_{3} = 8 - 1 = 7$, and at $\lambda = 1/6$, $\sum|B|^{2} = 31/18$ and $m^{2} = 31/9 + 7/3 = 52/9$.

The script checks, in exact Gaussian-rational arithmetic, the worked example, $\det\Omega_{\lambda}$ through $\det(I^{ab}) = \lambda^{2}$, the identity of (3) and the radical property for random triples (also the first three constituents of random four- and five-constituent composites) at four values of $\lambda$, nonzero faces and $m^{2} > 0$ for $\lambda > 0$ through the join, the $\Lambda = 0$ family of (4) with $m^{2} = 0$ and its massiveness once $\lambda \ne 0$, and the rank-two massless plane.
